\section{Dataset and Ground Truth}
\label{sec:dataset}

% Draft for Ammaar's revision. Facts sourced from data/README.md,
% docs/GT_VERIFICATION_GUIDE.md, and the split manifests — keep in sync.

\subsection{Source Dataset}
Digitize-HCD \cite{digitizehcd_data} (CC BY 4.0) contains 1{,}277
photographs of hand-drawn circuit diagrams at roughly 2000\,px
resolution, annotated with 18{,}600 COCO-format component bounding
boxes over 17 classes (BJT-NPN/PNP, capacitor, diode, Zener diode,
GND, AC/DC current and voltage sources, one-port DC source, inductor,
MOSFET-N/P, op-amp, resistor, and \emph{Wire Crossover}), MMOCR-style
text annotations with transcriptions, and per-class component crops
with port-location Gaussian heatmaps and pixel port coordinates. We
verified our copy byte-for-byte against the published archive (SHA-256
of the 2.1\,GB release; every one of the 1{,}277 images matched).

The COCO records carry no drafter identity, so drafter-disjoint
splitting is impossible — a stated limitation of any benchmark on this
dataset. We instead freeze a stratified split of 895/192/\NumTestImages{}
(train/val/test), stratified by component-count tertile and
rarest-class presence; all \NumClasses{} classes appear in every split,
with test-set support ranging from \SupportMin{}
(\SupportMinClass) to \SupportMax{} (\SupportMaxClass) instances
(Table~\ref{tab:detection}). The split manifests are versioned with the
code.

\subsection{Preprocessing and Coordinate Frames}
All stages operate on a \emph{cleaned frame}: deskew (length-weighted
median of Hough segment angles folded mod $90^\circ$, with a robust
fallback), shadow normalization, binarization, a content crop, and an
aspect-preserving resize onto a $512\times512$ canvas. Two properties
make this frame benchmark-safe. First, the full geometric transform of
every image is recorded, so published original-frame annotations
project exactly into the cleaned frame (verified: 0 of 18{,}600
projected boxes fall off-canvas, enforced by a guard script). Second,
the crop is \emph{annotation-aware}: every annotated component is
unioned into the crop rectangle before speck removal, making it
structurally impossible for a labeled component to be cropped out.

\subsection{Human-Verified Connectivity Ground Truth}
\label{sec:gt}
Digitize-HCD ships no connectivity ground truth, so we created it for
both evaluation splits. Each GT file is a component-terminal graph: per
component, a class, a bounding box, and per-terminal net assignments
(ground is always net~0), under a strict schema with a validator that
rejects unnamed terminals, single-terminal nets, and class/terminal
arity mismatches.

Critically, the test-split annotation was \emph{not} produced by running
the pipeline and correcting its output. Component inventory and box
geometry come from the published COCO annotations; every net was traced
from \texttt{null}. Each image went through wire extraction under
hysteresis binarization, a criticality analysis that flips every
junction/crossing call and retains only the sites where the partition of
terminals changes, adjudication by eye at 5--24$\times$ magnification,
deterministic net reconstruction from the recorded decisions, an
electrical rule check, and an automated self-consistency re-derivation.
Nets are recomputed from the decision record rather than hand-edited, so
every shipped netlist is reproducible from it, and the decision records
are released alongside the netlists.

Ground truth records the topology \emph{as drawn}. Where a faithful
reading yields a circuit that cannot be simulated --- one sheet in the
test split has a rail tip landing on a grounded column, which
short-circuits a voltage source --- the annotation keeps the drawn
topology and the electrical rule check reports the fault. Electrically
motivated corrections are never folded into the topology; the
observation that a shorted voltage source appears nowhere else in the
annotated set is recorded in that file's notes as evidence about the
drawing, not as a change applied to it.

That protocol matters for more than rigor: ground truth seeded by the
pipeline's own output would make any evaluation of that pipeline
partly circular. Table~\ref{tab:gt-verification} reports the volume of
judgement it took --- every count derived from the released decision
records, not from prose --- and Table~\ref{tab:splits} gives the split
composition. Two findings from the pass are worth carrying forward: the
dominant error mode is a terminal landing on a stroke of the handwritten
value label rather than on the lead, which is why a one-terminal net is a
hard error in the rule check; and pin order on three-terminal parts was
wrong far more often than net grouping, a failure no net-level or
electrical check can detect.

The re-derivation rows of Table~\ref{tab:gt-verification} must be read
for what they are. Each sampled file was re-derived from the drawing
without reading the first pass's reasoning first, and the two netlists
were compared --- but that re-derivation was performed by an \emph{AI
assistant, not by an independent human annotator}. It is a real check on
clerical and transcription error --- formatting faults, inconsistent
terminal counts, duplicate nodes, one-terminal nets --- and it caught
them: seven of the nine files flagged by the cross-checks changed. It is
a weak check on judgement, because it shares the first pass's reasoning
and blind spots, so agreement is evidence of consistency rather than of
correctness. It is not an inter-annotator agreement study and we do not
report it as one. Independent second annotation by an external annotator
is \emph{pending}.
\todoa{external second annotation of ambiguous crossings and
three-terminal pin order — state the result here if/when it happens.}
The verified GT files and decision records are released with the
benchmark; images remain under the original Digitize-HCD license.

\begin{table}[t]
  \centering
  \caption{Split composition. Connectivity ground truth exists for both
  evaluation splits; the training split carries only the published box
  annotations.}
  \label{tab:splits}
  % AUTO-GENERATED by scripts/make_paper_tables.py — do not edit.
\begin{tabular}{lrrrrl}
\toprule
split & images & components & terminals & nets & role \\
\midrule
train & 895 & -- & -- & -- & detector training \\
val & 190 & 2576 & 5116 & 1509 & parameter selection \\
test & 192 & 2564 & 5090 & 1496 & reported, never selected on \\
\bottomrule
\end{tabular}

\end{table}

\begin{table}[t]
  \centering
  \caption{Judgement recorded while building the test-split ground truth.
  Every count is derived from the released per-image decision records by
  \texttt{scripts/gt\_verification\_stats.py}. The lower block reports an
  \emph{automated self-consistency re-derivation}: files re-derived from
  the drawing before reading the first pass, by an AI assistant and
  \emph{not} by an independent human annotator. It detects clerical and
  transcription error; it cannot establish that a judgement call is
  correct, because it shares the first pass's reasoning. Independent
  human second annotation is pending.}
  \label{tab:gt-verification}
  % AUTO-GENERATED by scripts/make_paper_tables.py — do not edit.
\begin{tabular}{lr}
\toprule
recorded decision & count \\
\midrule
intersection sites adjudicated & 2047 \\
\quad read as junction & 1707 \\
\quad read as crossing & 194 \\
\quad explicit edge grouping & 63 \\
\quad ink meets, nothing joins & 83 \\
terminals repointed to a different lead & 1261 \\
nets asserted where the box swallowed the contact & 40 \\
wire fragments re-joined across a scan gap & 38 \\
gap bridges rejected as not-touching & 3 \\
published component classes corrected & 3 \\
components marked deliberately unconnected & 1 \\
\midrule
\multicolumn{2}{l}{\emph{second reader, re-derived from the drawing first}} \\
\quad blind random & 0/12 \\
\quad stratified on sheets with 3-terminal parts (38 3-terminal parts) & 0/5 \\
\quad flagged by cross-checks & 7/9 \\
\bottomrule
\end{tabular}

\end{table}

\subsection{Cross-Dataset Material}
For context and future zero-shot evaluation we also stage CGHD
\cite{cghd} (3{,}255 images, 25 drafters) with a mapping from its
53-entry class dictionary onto the 17 Digitize-HCD categories.
\draftnote{keep only if the zero-shot experiment [IDEAL] actually runs;
otherwise fold into future work.}
